% Vitali选择的放大论证：已选不交区间 I_1=[0.5,2.5]（长 2）、I_2=[3.4,4.6]（长 1.2）、I_3=[5.6,6.3]（长 0.7）；未选区间 I=[2.2,3.0]（长 0.8）与 I_1 相交，且 |I|<2|I_1|，所以 I 落在与 I_1 同心、长度 5 倍的区间 [-3.5,6.5] 内（图中截取到 [-0.4,6.5] 部分）。
\begin{tikzpicture}[>=Stealth,
  sel/.style={draw=FigureBlue,fill=FigureBlueFill,line width=.8pt},
  cand/.style={draw=FigureRed,fill=FigureRedFill,line width=.8pt},
  note/.style={font=\FigureNoteFont\fontsize{7.5}{9}\selectfont,text=FigureMuted}]
\draw[FigureStroke,line width=.6pt,->] (-.5,0)--(7.0,0);
\foreach \a/\b/\n in {0.5/2.5/1,3.4/4.6/2,5.6/6.3/3}{
  \draw[sel] (\a,.15) rectangle (\b,.45);
  \node at ({(\a+\b)/2},.3) {$I_{\n}$};}
\draw[cand] (2.2,.7) rectangle (3.0,1.0);
\node at (2.6,.85) {$I$};
\fill[FigureStroke] (2.75,0) circle (.05) node[below] {$x$};
\draw[FigureStroke,dashed,line width=.4pt] (2.75,0)--(2.75,.7);
\draw[FigureGreen,line width=.9pt,dashed] (-.4,-.55)--(6.5,-.55);
\draw[FigureGreen,line width=.9pt] (6.5,-.45)--(6.5,-.65);
\node[note,anchor=west] at (-.4,-.8) {$5I_1$：与$I_1$同心、长度为$5|I_1|$（左端在图外）};
\end{tikzpicture}
